\section{附录：核心变量与标注规则摘要}

\subsection{核心变量定义}

\begin{longtable} {p{3cm}p{4cm}p{7cm}}
    \toprule
    变量名            & 中文名称         & 度量说明                                        \\
    \midrule
    \endfirsthead
    \toprule
    变量名            & 中文名称         & 度量说明                                        \\
    \midrule
    \endhead
    AuthIndex      & AI 披露真实性综合指数 & 对提及强度、场景具体性和表述审慎性标准化后按 $0.20/0.45/0.35$ 加权。 \\
    PackagingIndex & 叙事包装性指数      & 基于包装性评分及相关包装特征形成的反向指标。                      \\
    ROA            & 上市后 ROA      & 上市后第一年净利润与平均总资产之比。                          \\
    RevenueGrowth  & 营业收入增长率      & 上市后第一年营业收入同比增速。                             \\
    AssetTurnover  & 总资产周转率       & 上市后第一年营业收入与平均总资产之比。                         \\
    PerfIndex      & 综合经营表现指数     & 对 ROA、营业收入增长率和总资产周转率标准化后取均值。                \\
    LnAssets       & 企业规模         & 总资产取自然对数。                                   \\
    Leverage       & 资产负债率        & 总负债除以总资产。                                   \\
    FirmAge        & 企业年龄         & 上市年份减去成立年份。                                 \\
    RDIntensity    & 研发强度         & 研发费用除以营业收入。                                 \\
    SOE            & 所有制性质        & 国有企业取 1，否则取 0。                              \\
    \bottomrule
\end{longtable}

\subsection{标注规则摘要}

\begin{table}[H]
    \centering
    \caption{文本标注维度摘要}
    \begin{tabular}{p{3cm}p{9.5cm}}
        \toprule
        维度    & 核心判断标准                               \\
        \midrule
        提及强度  & 看关键词频次、相关段落数和文本占比，识别企业“说了多少”。        \\
        场景具体性 & 看是否落到具体业务场景、流程节点、系统模块、组织安排和应用结果。     \\
        表述审慎性 & 看是否说明实施边界、阶段进度和限制条件，是否区分“已落地”和“拟推进”。 \\
        包装性   & 看是否存在宏大口号、概念堆砌、承诺化措辞和明显夸大修辞。         \\
        \bottomrule
    \end{tabular}
\end{table}
